\chapter{All-Loop Asymptotic Bethe Ansatz}
\label{ch:planar-n4-asymptotic-bethe-ansatz}

\ConventionsLorentzian

The one-loop spin chain of
Chapter~\ref{ch:planar-n4-spectral-problem-spin-chains} is only the first
term in the planar dilatation operator.  The all-loop asymptotic Bethe ansatz
is a proposal for the spectrum of long single-trace operators before wrapping
effects become important.  Its inputs are the residual centrally extended
\(\mathfrak{su}(2|2)\oplus\mathfrak{su}(2|2)\) symmetry of the BMN vacuum,
factorized magnon scattering, crossing, and analyticity.

\section{BMN Vacuum and Central Extension}

The vacuum spin-chain state is
\[
  \operatorname{tr}(Z^J).
\]
It is protected and carries large \(R\)-charge \(J\).  A magnon is an
insertion into this long chain.  The residual symmetry preserving the vacuum
contains two centrally extended \(\mathfrak{su}(2|2)\) algebras.  For one
copy, write the bosonic states as \(\ket{\phi^a}\), \(a=1,2\), and the
fermionic states as \(\ket{\psi^\alpha}\), \(\alpha=1,2\).  The supercharges
act in the form
\[
\begin{aligned}
  Q^\alpha{}_a\ket{\phi^b}
  &=a\,\delta_a^b\ket{\psi^\alpha},
&
  Q^\alpha{}_a\ket{\psi^\beta}
  &=b\,\varepsilon^{\alpha\beta}\varepsilon_{ab}\ket{\phi^b Z^+},
\\
  S^a{}_\alpha\ket{\phi^b}
  &=c\,\varepsilon^{ab}\varepsilon_{\alpha\beta}\ket{\psi^\beta Z^-},
&
  S^a{}_\alpha\ket{\psi^\beta}
  &=d\,\delta_\alpha^\beta\ket{\phi^a}.
\end{aligned}
\]
The symbols \(Z^+\) and \(Z^-\) record insertion or removal of one vacuum
letter.  On a momentum eigenstate,
\[
  Z^+\mapsto \ee^{\ii p}-1,
  \qquad
  Z^-\mapsto \ee^{-\ii p}-1
\]
up to the chosen side convention.  The representation constraint is
\[
  ad-bc=1.
\]
The central charges are
\[
  P=ab,\qquad K=cd,\qquad H=ad+bc .
\]
The shortening condition of the centrally extended algebra is
\[
  H^2-4PK=1 .
\]
With
\[
  P=g(1-\ee^{\ii p}),\qquad K=g(1-\ee^{-\ii p}),
\]
this gives the exact single-magnon dispersion relation
\begin{equation}
  H(p)=\sqrt{1+16g^2\sin^2\frac p2}.
  \label{eq:planar-n4-magnon-dispersion}
\end{equation}
The anomalous contribution of a magnon is \(H(p)-1\).

\begin{proposition}[Dispersion from shortening]
\label{prop:planar-n4-dispersion-shortening}
Assuming the centrally extended residual symmetry acts on a one-magnon
representation with central charges \(P=g(1-\ee^{\ii p})\) and
\(K=g(1-\ee^{-\ii p})\), the shortening condition fixes the dispersion
relation to \eqref{eq:planar-n4-magnon-dispersion}.
\end{proposition}

\begin{proof}
Compute
\[
  PK
  =
  g^2(1-\ee^{\ii p})(1-\ee^{-\ii p})
  =
  4g^2\sin^2\frac p2 .
\]
Substitution into \(H^2-4PK=1\) gives
\[
  H^2=1+16g^2\sin^2\frac p2 .
\]
The positive square root is selected because \(H=1\) at \(g=0\) for an
elementary magnon.
\end{proof}

\section{Zhukovsky Variables}

Introduce spectral parameters \(x^\pm\) by
\begin{equation}
  \ee^{\ii p}=\frac{x^+}{x^-},
  \qquad
  x^++\frac1{x^+}-x^- -\frac1{x^-}=\frac{\ii}{g}.
  \label{eq:planar-n4-xpm-definition}
\end{equation}
The physical region for real magnons is chosen so that
\((x^+)^*=x^-\) and \(|x^\pm|>1\) at weak coupling.  The magnon energy is
\[
  H
  =
  1+\ii g\left(
    \frac1{x^+}-\frac1{x^-}
  \right)
  -
  \ii g(x^+-x^-),
\]
with the same branch as \eqref{eq:planar-n4-magnon-dispersion}.  Conserved
charges used in the dressing phase are
\[
  q_r(x^\pm)
  =
  \frac{\ii}{r-1}
  \left[
    \frac1{(x^+)^{r-1}}-\frac1{(x^-)^{r-1}}
  \right],
  \qquad r\geq2 .
\]

\section{Two-Body S-Matrix}

The magnon two-body \(S\)-matrix is fixed by symmetry up to a scalar factor.
Let \(S^{\rm mat}_{12}\) denote the matrix part determined by the two
centrally extended \(\mathfrak{su}(2|2)\) copies.  The full two-magnon
operator is
\[
  S_{12}=S_0(x_1^\pm,x_2^\pm)\,S^{\rm mat}_{12}.
\]
The scalar factor is written as
\[
  S_0(x_1^\pm,x_2^\pm)
  =
  S_{\rm rat}(x_1^\pm,x_2^\pm)\,
  \sigma(x_1^\pm,x_2^\pm)^2,
  \qquad
  \sigma=\ee^{\ii\theta}.
\]
The rational part is fixed by the same representation data that determine
the matrix part.  The dressing phase has the antisymmetric charge expansion
\[
  \theta(x_1^\pm,x_2^\pm)
  =
  \sum_{r=2}^\infty\sum_{\substack{s>r\\ r+s\ {\rm odd}}}
  c_{r,s}(g)\,
  \bigl(q_r(x_1^\pm)q_s(x_2^\pm)-q_s(x_1^\pm)q_r(x_2^\pm)\bigr).
\]
Crossing relates \(\sigma(x_1^\pm,x_2^\pm)\) to the phase obtained by
analytically continuing the first magnon through the crossed sheet
\(x_1^\pm\mapsto1/x_1^\pm\).  It supplies an additional analyticity condition
on the scalar factor beyond factorized Yang--Baxter compatibility.

\begin{quotedtheorem}[Asymptotic planar magnon S-matrix]
\label{qthm:planar-n4-asymptotic-smatrix-status}
Assume planar integrability of the long \(\mathcal N=4\) SYM spin chain, the
centrally extended residual symmetry, physical unitarity, Yang--Baxter
factorization, crossing, and the minimal analyticity condition used in the
BES construction.  Then the asymptotic two-body magnon S-matrix is the
\(\mathfrak{su}(2|2)^2\)-invariant matrix part multiplied by the scalar
dressing factor whose phase has the charge expansion above with the BES
coefficients \(c_{r,s}(g)\).
\end{quotedtheorem}

This is stated as a quoted theorem because the full claim is not a theorem of
four-dimensional QFT derived from a nonperturbative definition of the gauge
theory.  It is an integrability theorem inside a framework whose physical
identification with the planar gauge-theory dilatation operator has extensive
checks and strong structural constraints.

\section{Asymptotic Bethe Equations}

In the \(SU(2)\) scalar subsector, the all-loop asymptotic Bethe equations can
be written
\begin{equation}
  \left(\frac{x_j^+}{x_j^-}\right)^L
  =
  \prod_{\substack{k=1\\ k\neq j}}^M
  \frac{x_j^- -x_k^+}{x_j^+ -x_k^-}
  \frac{1-1/(x_j^+x_k^-)}{1-1/(x_j^-x_k^+)}
  \sigma(x_j^\pm,x_k^\pm)^2 ,
  \label{eq:planar-n4-su2-aba}
\end{equation}
with cyclicity
\[
  \prod_{j=1}^M\frac{x_j^+}{x_j^-}=1.
\]
The scaling dimension is
\[
  \Delta
  =
  L+\sum_{j=1}^M\left(
    \sqrt{1+16g^2\sin^2\frac{p_j}{2}}-1
  \right),
\]
up to wrapping corrections.  The full
\(\mathfrak{psu}(2,2|4)\) asymptotic Bethe equations require nested roots
for the residual symmetry algebra.  Their algebraic role is the same as in
Chapter~\ref{ch:integrable-nested-bethe-ansatz-matrix-bethe-yang}; their
physical meaning is different because these auxiliary roots organize
polarizations of magnons on a gauge-theory trace.

\begin{controlledapproximation}[Asymptotic long-chain regime]
\label{ca:planar-n4-asymptotic-regime}
Equation~\eqref{eq:planar-n4-su2-aba} applies to the planar single-trace
spectrum in the regime where the interaction range is shorter than the trace
length \(L\).  At weak coupling, wrapping effects begin at the loop order at
which planar interactions can close around the trace.  At finite coupling,
the distinction is between the asymptotic Bethe-Yang problem and the
finite-size problem solved by mirror TBA or the quantum spectral curve.
\end{controlledapproximation}

\section{Weak-Coupling Reduction}

At small \(g\), \(x^\pm\) reduce to the one-loop rapidity variable:
\[
  x^\pm=\frac{1}{g}\left(u\pm\frac{\ii}{2}\right)+O(g).
\]
Then
\[
  \frac{x_j^+}{x_j^-}
  =
  \frac{u_j+\ii/2}{u_j-\ii/2}+O(g^2),
\]
and the rational part of \eqref{eq:planar-n4-su2-aba} reduces to
\[
  \frac{u_j-u_k+\ii}{u_j-u_k-\ii}+O(g^2).
\]
The dressing phase starts at higher weak-coupling order, so the leading
equation is the one-loop Bethe equation
\eqref{eq:planar-n4-one-loop-su2-bethe}.  The dispersion expands as
\[
  H(p)-1
  =
  8g^2\sin^2\frac p2
  -32g^4\sin^4\frac p2
  +256g^6\sin^6\frac p2
  +O(g^8).
\]
Using \(g^2=\lambda/(16\pi^2)\), the leading term agrees with the one-loop
normalization of Chapter~\ref{ch:planar-n4-spectral-problem-spin-chains}.
